% GENERATED FILE. Edit canonical lessons, then run npm run generate:books.
% canonical-source-hash: fnv1a64:0844a76325bb31c2
% canonical-lessons: LA-C450-querela, LA-C450-conqueror, LA-C450-fraudo, LA-C450-rapax, LA-C450-sordes

\chapter{Complaint, To complain, To cheat, Grasping, Dirt}
\label{ch:la-complaint-to-complain-to-cheat-grasping-dirt}
\hlchaptermodality{450}

\begin{chapteropening}
\textbf{By the end of this chapter:} \emph{I can complain about a bad inn: making a complaint, complaining, being cheated, a grasping landlord, and dirt.}

Complete the last lesson of chapter 450: I can complain about a bad inn: making a complaint, complaining, being cheated, a grasping landlord, and dirt.
\end{chapteropening}

\section[querēla, querēlae]{querēla, querēlae — a complaint}

\label{lesson:LA-C450-querela}



\begin{quote}
Before the new one: say the Latin for rotten, then the Latin for a tavern hostess.
\end{quote}

\subsection*{You'll want to know: querēla, querēlae}
\textbf{querēla, querēlae} — \textquotedblleft{}a complaint\textquotedblright{}.

\begin{quote}
Hospitēs multās querēlās dē caupōne habent.
\end{quote}

\begin{quote}
\emph{The guests have many complaints about the innkeeper.}
\end{quote}

First declension, feminine: \emph{querēla, querēlae f.}, also spelled \emph{querella}. English \emph{quarrel} comes from it through French.

\begin{cousinweb}
\emph{querēla} is built on \emph{queror} 'to complain' with the suffix \emph{-ēla} (EtymDB; Lewis \& Short). English \emph{quarrel} comes from it through Old French.
\end{cousinweb}

\begin{grammarlens}[title={What you've built}]
One more word for telling, explaining and describing at B1.
\end{grammarlens}

\subsection*{Your turn}
\begin{itemize}
\raggedright
  \item \emph{Say it:} \emph{querēla, querēlae}
  \item \emph{Say it:} \emph{querēla, querēlae}, once more
  \item \emph{Say it:} a sentence of your own with \emph{querēla, querēlae}
  \item \emph{Recall:} say the Latin for rotten, then the Latin for a tavern hostess, then say \emph{querēla, querēlae} again
\end{itemize}

\begin{tcolorbox}[breakable,colback=teal!4,colframe=teal!35!black,title={Before you move on}]
What does querēla, querēlae mean? (\textquotedblleft{}A complaint\textquotedblright{}.) Say it once more.
\end{tcolorbox}
\section[conqueror, conquerī]{conqueror, conquerī — to complain (loudly), to complain of}

\label{lesson:LA-C450-conqueror}



\begin{quote}
Before the new one: say the Latin for a tavern hostess, then the Latin for a complaint.
\end{quote}

\subsection*{You'll want to know: conqueror, conquerī}
\textbf{conqueror, conquerī} — \textquotedblleft{}to complain (loudly), to complain of\textquotedblright{}.

\begin{quote}
Viātor dē cīmicibus et dē pretiō conqueritur.
\end{quote}

\begin{quote}
\emph{The traveller complains about the bedbugs and about the price.}
\end{quote}

Deponent: \emph{conqueror, conquerī, conquestus sum}. English \emph{conquer} does not come from it; that word goes back to \emph{conquīrere} 'to seek out'.

\begin{cousinweb}
\emph{conqueror} is built from \emph{con-}, which strengthens, and \emph{queror} 'to complain' (EtymDB; Lewis \& Short). English \emph{conquer} looks like it but comes from a different verb, \emph{conquīrere} 'to seek out'.
\end{cousinweb}

\begin{grammarlens}[title={What you've built}]
One more word for telling, explaining and describing at B1.
\end{grammarlens}

\subsection*{Your turn}
\begin{itemize}
\raggedright
  \item \emph{Say it:} \emph{conqueror, conquerī}
  \item \emph{Say it:} \emph{conqueror, conquerī}, once more
  \item \emph{Say it:} a sentence of your own with \emph{conqueror, conquerī}
  \item \emph{Recall:} say the Latin for a tavern hostess, then the Latin for a complaint, then say \emph{conqueror, conquerī} again
\end{itemize}

\begin{tcolorbox}[breakable,colback=teal!4,colframe=teal!35!black,title={Before you move on}]
What does conqueror, conquerī mean? (\textquotedblleft{}To complain (loudly), to complain of\textquotedblright{}.) Say it once more.
\end{tcolorbox}
\section[fraudō, fraudāre]{fraudō, fraudāre — to cheat, to defraud}

\label{lesson:LA-C450-fraudo}



\begin{quote}
Before the new one: say the Latin for a complaint, then the Latin for to complain.
\end{quote}

\subsection*{You'll want to know: fraudō, fraudāre}
\textbf{fraudō, fraudāre} — \textquotedblleft{}to cheat, to defraud\textquotedblright{}.

\begin{quote}
Caupō malus hospitēs fraudat.
\end{quote}

\begin{quote}
\emph{The bad innkeeper cheats his guests.}
\end{quote}

\emph{fraudō, fraudāre, fraudāvī, fraudātum}, with the person in the accusative and the thing in the ablative. English \emph{defraud} comes from \emph{dēfraudāre}.

\begin{cousinweb}
\emph{fraudō} is built on \emph{fraus, fraudis} 'harm, deceit' (EtymDB; Lewis \& Short). English \emph{fraud} comes from the noun and \emph{defraud} from the compound \emph{dēfraudāre}.
\end{cousinweb}

\begin{grammarlens}[title={What you've built}]
One more word for telling, explaining and describing at B1.
\end{grammarlens}

\subsection*{Your turn}
\begin{itemize}
\raggedright
  \item \emph{Say it:} \emph{fraudō, fraudāre}
  \item \emph{Say it:} \emph{fraudō, fraudāre}, once more
  \item \emph{Say it:} a sentence of your own with \emph{fraudō, fraudāre}
  \item \emph{Recall:} say the Latin for a complaint, then the Latin for to complain, then say \emph{fraudō, fraudāre} again
\end{itemize}

\begin{tcolorbox}[breakable,colback=teal!4,colframe=teal!35!black,title={Before you move on}]
What does fraudō, fraudāre mean? (\textquotedblleft{}To cheat, to defraud\textquotedblright{}.) Say it once more.
\end{tcolorbox}
\section[rapāx, rapācis]{rapāx, rapācis — grasping, greedy}

\label{lesson:LA-C450-rapax}



\begin{quote}
Before the new one: say the Latin for to complain, then the Latin for to cheat.
\end{quote}

\subsection*{You'll want to know: rapāx, rapācis}
\textbf{rapāx, rapācis} — \textquotedblleft{}grasping, greedy\textquotedblright{}.

\begin{quote}
Caupō rapāx magnum pretium prō cēnā parvā rogat.
\end{quote}

\begin{quote}
\emph{The grasping innkeeper asks a big price for a small dinner.}
\end{quote}

An adjective of one ending: \emph{rapāx}, genitive \emph{rapācis}. Also of wild beasts: ravenous. English \emph{rapacious} comes from it.

\begin{cousinweb}
\emph{rapāx} is built on \emph{rapiō} 'to seize, to snatch' with the suffix \emph{-āx} 'given to' (EtymDB; Lewis \& Short). English \emph{rapacious} comes from it.
\end{cousinweb}

\begin{grammarlens}[title={What you've built}]
One more word for telling, explaining and describing at B1.
\end{grammarlens}

\subsection*{Your turn}
\begin{itemize}
\raggedright
  \item \emph{Say it:} \emph{rapāx, rapācis}
  \item \emph{Say it:} \emph{rapāx, rapācis}, once more
  \item \emph{Say it:} a sentence of your own with \emph{rapāx, rapācis}
  \item \emph{Recall:} say the Latin for to complain, then the Latin for to cheat, then say \emph{rapāx, rapācis} again
\end{itemize}

\begin{tcolorbox}[breakable,colback=teal!4,colframe=teal!35!black,title={Before you move on}]
What does rapāx, rapācis mean? (\textquotedblleft{}Grasping, greedy\textquotedblright{}.) Say it once more.
\end{tcolorbox}
\section[sordēs, sordis]{sordēs, sordis — dirt, filth}

\label{lesson:LA-C450-sordes}



\begin{quote}
Before the new one: say the Latin for to cheat, then the Latin for grasping.
\end{quote}

\subsection*{You'll want to know: sordēs, sordis}
\textbf{sordēs, sordis} — \textquotedblleft{}dirt, filth\textquotedblright{}.

\begin{quote}
In cubiculō sordēs sunt, et hospes īrātus est.
\end{quote}

\begin{quote}
\emph{There is dirt in the bedroom, and the guest is angry.}
\end{quote}

Third declension, feminine: \emph{sordēs, sordis f.}, mostly plural \emph{sordēs, sordium}. English \emph{sordid} comes from the adjective \emph{sordidus}.

\begin{cousinweb}
\emph{sordēs} is built on \emph{sordeō} 'to be dirty' (EtymDB; Lewis \& Short). English \emph{sordid} comes from the related adjective \emph{sordidus}.
\end{cousinweb}

\begin{grammarlens}[title={What you've built}]
One more word for telling, explaining and describing at B1.
\end{grammarlens}

\subsection*{Your turn}
\begin{itemize}
\raggedright
  \item \emph{Say it:} \emph{sordēs, sordis}
  \item \emph{Say it:} \emph{sordēs, sordis}, once more
  \item \emph{Say it:} a sentence of your own with \emph{sordēs, sordis}
  \item \emph{Recall:} say the Latin for to cheat, then the Latin for grasping, then say \emph{sordēs, sordis} again
\end{itemize}

\begin{tcolorbox}[breakable,colback=teal!4,colframe=teal!35!black,title={Before you move on}]
What does sordēs, sordis mean? (\textquotedblleft{}Dirt, filth\textquotedblright{}.) Say it once more.
\end{tcolorbox}
